The design was frozen at a tagged commit, given in Section~\ref{sec:repro}, before
any measurement. Six amendments were appended after the freeze. Each is a dated
section of the pre-registration document that states what changed and what had been
seen when it was written. Nothing above an amendment was edited, and no amendment
deletes text.

\begin{description}
\item[Amendment 1 (2026-10-20).] Corrected the seed claim: one SplitMix64
step, not two, and the formatting rules that make the seed message round-trip.
Written after the freeze and before any measurement.
\item[Amendment 2 (2026-10-20).] Corrected the anchor's relation to the earlier
recovery study: equal to within $8\times10^{-5}$ relative, not exactly equal.
Written after the freeze and before any measurement.
\item[Amendment 3 (2026-10-20).] Fixed the seed reading, the burn-in rule and the
resolved-cell rule for the shallow-decay half of the grid. Written after a
twenty-replication pilot in eight cells and before any full-grid cell ran; no
hypothesis was evaluated on the pilot.
\item[Amendment 4 (2026-10-20).] Added robust companions, which decide nothing, and
restricted the headline's pooling to the two-type cells. Written after the full grid
had run and a counts-only pass over the raw rows had been read; that pass showed the
heavy tail in the decay estimate, and the amendment says so.
\item[Amendment 5 (2026-10-20).] Fixed the Holm family of four, the $z$ and $p$ for
each rule, the bootstrap standard error of $q$, and the population of the
quantisation hypothesis. Written after the counts-only pass and before any estimand
was computed.
\item[Amendment 6 (2026-10-21).] Resolved the quantisation hypothesis's cell
population and adopted the rule that where admissible forms disagree, the least
favourable label governs. Written after the analysis, before the paper.
\end{description}

One command per stage, each run from the repository root:

\begin{description}
\item[the grid] \texttt{python -m studies.run\_study}
\item[the estimands] \texttt{python -m studies.analyze} and
  \texttt{python -m studies.phase3}
\item[the shipped map] \texttt{python -m studies.build\_map}
\item[the figures] \texttt{python -m studies.figures.make\_figures}
\item[the interval study] \texttt{python -m studies.run\_auto} and
  \texttt{python -m studies.analyze\_auto}
\item[this document] \texttt{python -m studies.make\_paper\_tex}
\end{description}
